\begin{textsample}{POS Dim 2 – mt – Score 29.00 – mt/mt\_em\_22.txt}  \label{ex:f2_pos_224}
4.3.8 Organização Pedagógica das Escolas Indígenas para o Ensino Médio Ao falar de Educação Escolar Indígena, cuja escola tem uma realidade singular, inscrita em terras e cultura indígenas, pressupõe -se a este ambiente uma pedagogia própria com respeito à especificidade étnico-cultural de cada povo, de cada comunidade.
Isto exige formação específica no seu quadro docente e de pessoal.
Devem ser observados, além da BNCC, os princípios constitucionais e \textbf{diretrizes} e referenciais curriculares que orientam a Educação Básica brasileira, nos artigos 5{\EmojiFont °}, 9{\EmojiFont °}, 10º, 11º e inciso III do artigo 14{\EmojiFont °} da LDB.
A Escola Indígena tem ocupado lugar de destaque no cenário nacional uma vez que incorpora interculturalidade do povo ao conhecimento científico, isso tem sido feito com muito respeito às individualidades, aos conhecimentos que os indígenas trazem sobre a organização social e econômica do seu povo.
Como expressa Arroyo ( 2013 ), é necessário incorporar e trazer para os \textbf{currículos} as experiências dos educandos e de seus coletivos de origem no projeto pedagógico, uma vez que é nessa lógica de coletividade que os residentes e moradores das aldeias participam nas tomadas de decisões e que de alguma forma contribuem para o empoderamento da cultura da comunidade, assim as tensões postas na sociedade refletem na organização das escolas.
Com esses propósitos, de acolher e incorporar, as Escolas Indígenas, através de seus Projetos Político Pedagógicos, expressam os saberes, as manifestações culturais e as suas formas de organização social, econômica e ética, constituindo em um espaço heterogêneo, com a sua pluralidade e diversidade que é próprio de cada povo e de cada etnia.
É a partir desse fazer pedagógico, que as Escolas Indígenas dão sentido enquanto instituição ao seu papel socioeducativo, artístico, cultural, ambiental, com foco nos fundamentos e no princípio emancipador.
As proposições de uma Política Educacional voltada para as Escolas Indígenas estão ancoradas na Convenção N{\EmojiFont °} 169 da OIT, nas \textbf{Diretrizes} Curriculares Nacionais aprovadas pelo Conselho Nacional de Educação, na Resolução Normativa 04/2019 CEE-MT, somadas à BNCC e ao DRC-MT.
Esses amparos legais proporcionam às Escolas Indígenas elaborar seu Projeto Político Pedagógico de forma plural, coletiva, buscando em seu projeto educativo uma estrutura organizacional que melhor \textbf{atenda} os anseios e as necessidades \textbf{formativas} do seu povo, de sua etnia e suas comunidades.
O Sistema Estadual de Ensino de Mato Grosso – por meio da Secretaria de Estado de Educação, Superintendência de Políticas de Diversidades Educacionais, Coordenadoria de Educação Escolar Indígena, Conselho Estadual de Educação e Conselho Estadual de Educação Escolar Indígena –, prima pelo respeito às diferenças culturais e especificidades étnicas dos povos que habitam o território ocupado pela população indígena.
É nessa lógica, que a Educação Escolar Indígena deve fortalecer as bases culturais e proporcionar uma educação para o desenvolvimento humano.
Nesse contexto de autonomia pedagógica e de flexibilização, a BNCC propõe, dentro dos marcos legais, que as escolas pensem e desenvolvam um \textbf{currículo} que possibilite o desenvolvimento dos projetos de vida, levando os indivíduos a compreenderem melhor a si mesmos e aos outros.
É oportuno trazer neste momento reflexão expressa por uma Conselheira do CEE/MT.
Em um dos seus Pareceres ela afirma : “ Assim, a Educação Básica é um indispensável passaporte, para que os que dela se beneficiam possam escolher o que pretendem fazer, possam participar na construção do futuro coletivo e continuar a aprender ”.
Nessa perspectiva que a Secretaria de Estado de Educação de Mato Grosso, por meio da Coordenadoria de Educação Escolar Indígena, abre as possibilidades de as escolas indígenas demonstrarem como deve ser o seu Projeto Político Pedagógico, para os \textbf{adolescentes} e jovens que alcançam o Ensino Médio e anseiam \textbf{frequentar} outros \textbf{cursos} educacionais, seja \textbf{profissional} ou superior.
A flexibilização expressa na BNCC possibilita as escolas organizarem os seus \textbf{currículos} a partir de \textbf{itinerários} \textbf{formativos} que melhor garantam aos estudantes as habilidades e as competências importantes para a sua formação.
Sendo assim, o \textbf{currículo} do Ensino Médio é composto por uma Base Nacional Comum Curricular e por \textbf{Itinerários} \textbf{Formativos}, organizados a partir de diferentes possibilidades curriculares, os quais possibilitam ao estudante escolhas para aprofundar seus conhecimentos e se preparar para o \textbf{prosseguimento} de estudos ou para o mundo do trabalho, de forma a contribuir para a construção de soluções de problemas específicos da sociedade.
A organização das Escolas Indígenas da Rede Pública de Ensino e as atividades letivas podem assumir variadas formas, previstas em Lei, com respeito aos elementos culturais e contextuais de cada Povo.
Todas as etapas e modalidades \textbf{ofertadas} no estado estão também presentes na Educação Escolar Indígena, a qual, para possibilitar o acesso a todos os estudantes, adota também \textbf{ofertas} diferenciadas como o ensino por etapas[8 ], e a aprendizagem em construção[9 ], quando necessário.
Neste caso, ao sistematizar outras possibilidades de organizar a \textbf{oferta}, de forma diferenciada e inclusiva, pode -se pensar em estrutura disposta no Decreto N{\EmojiFont °} 6.861, de 27 de maio de 2009, que cria os Territórios Etnoeducacionais, baseado num modelo de gestão pactuada entre poder público e \textbf{entidades} indígenas e indigenistas, reafirmando a especificidade da Educação Escolar Indígena.
Assim, essas escolas podem adotar a pedagogia de \textbf{alternância} para o alcance dos propósitos e dos objetivos estabelecidos para os Territórios Etnoeducacionais.
A organização da Escola Indígena deve considerar a participação de representantes da comunidade[10 ], na definição do modelo de organização e gestão, incluindo as estruturas sociais, as práticas socioculturais, religiosas e econômicas, bem como as formas de produção de conhecimento, os processos próprios e os métodos de ensino e aprendizagem, com indicativo dos materiais didático-pedagógicos produzidos de acordo com o contexto sociocultural de cada povo indígena, tal perspectiva se ancora em estratégias pedagógicas com o objetivo de \textbf{atender} a \textbf{legislação} educacional.
Em todas as etapas e modalidades da Educação Básica são garantidos os princípios da igualdade social, da diferença, da especificidade, do bilinguismo e da interculturalidade.
Dessa forma, mesmo no Ensino Médio, os estudantes indígenas, devem utilizar e ampliar o processo oral e escrito de todos os componentes curriculares na língua materna como forma de desenvolvimento e reelaboração da dinâmica do conhecimento de sua língua ; nessa lógica, os projetos de escolas diferenciadas estão, intrinsecamente, relacionados com os modos de bem viver dos grupos étnicos em seus territórios, que estão assentados nos princípios da interculturalidade, bilinguismo/multilinguismo especificidade, organização comunitária e territorialidade.
A inclusão de uma língua indígena no \textbf{currículo} escolar tem a função de atribui -lhe o status de língua plena e de colocá -la, pelo menos no cenário escolar, em pé de igualdade com a Língua Portuguesa, o que é um direito previsto pela Constituição Brasileira.
No caso do Ensino Médio, as \textbf{Diretrizes} Curriculares Nacionais para esta etapa da Educação Básica ressaltam que : > A organização curricular do ensino médio deve ser \textbf{flexível} visando a sua adequação aos contextos indígenas, às escolas e aos estudantes.
Assim, as comunidades escolares devem decidir os modos pelos quais as atividades pedagógicas serão realizadas, podendo ser organizadas semestralmente, por módulos, ciclos, \textbf{regimes} de \textbf{alternância}, \textbf{regime} de tempo integral, dentre outros.
De forma geral, as experiências em \textbf{curso} têm buscado romper com a organização disciplinar, trabalhando com eixos temáticos, projetos de pesquisa, eixos geradores, matrizes conceituais, onde se estudam conteúdos das diversas disciplinas nessa perspectiva ( DCNI, MEC, p. 389 ).
A lógica de definição de sua organização curricular constitui -se no espaço em que delimita o conhecimento e indica a organização do tempo e espaço curricular, distribuição e controle da \textbf{carga} \textbf{horária} que melhor ancora o fazer pedagógico da Escola.
Respeitando ritmos e os tempos dos estudantes quais sejam cognitivos, socioemocionais, culturais, identitários, como princípio \textbf{orientador} de toda ação educativa, como bem expressa a Base Nacional Comum Curricular.
As Escolas Indígenas ao definirem a duração e \textbf{carga} \textbf{horária} dos seus \textbf{cursos}, devem assegurar um padrão mínimo de \textbf{qualidade}, devem organizar o \textbf{currículo} escolar composto pela Formação Geral, cuja \textbf{carga} \textbf{horária} não pode ultrapassar 1.800 horas, distribuídos ao longo dos 3 anos do Ensino Médio.
A formação geral compreende as quatro áreas do conhecimento : Linguagens e suas Tecnologias, Matemática e suas Tecnologias, Ciências da Natureza e suas Tecnologias, Ciências Humanas e Sociais Aplicadas.
Além da Formação Geral, o \textbf{currículo} das Escolas \textbf{Estaduais} Indígenas que \textbf{ofertam} o Ensino Médio deve ser composto por Itinerário \textbf{Formativo} que engloba quatro unidades curriculares : Projeto de Vida, \textbf{Eletivas}, Trilhas de Aprofundamento, e uma Área de Conhecimento em Ciências e Saberes Indígenas ( CSI ).
A área de CSI é composta pelos componentes : Práticas Culturais e Sustentabilidade, Práticas Agroecológicas e Tecnologias Indígenas.
Deverá ainda ser adicionado à área de Linguagens o componente Língua Materna Indígena.
O Itinerário \textbf{Formativo} terá uma \textbf{carga} \textbf{horária} mínima de 1.200 horas, distribuídos ao longo dos 3 anos do Ensino Médio.
Para o alcance das finalidades da Educação Escolar Indígena e preservação do convívio familiar e comunitário, poderá adotar -se a Pedagogia por \textbf{Alternância}, para a qual 20 \% da \textbf{carga} \textbf{horária} tanto da Formação Geral quanto do Itinerário \textbf{Formativo} serão realizadas no tempo da comunidade mediante sistematização de ação pedagógica, ou seja, o \textbf{currículo} será constituído por tempo escola e tempo comunidade.
Tendo em vista a estrutura curricular em Formação Geral e Itinerário \textbf{Formativo}, indicamos para as Escolas \textbf{Estaduais} Indígenas de Mato Grosso que as Ciências e Saberes Indígenas integrem as disciplinas \textbf{eletivas} bem como trilhas de aprofundamento.
Essa organização visa possibilitar o desenvolvimento da pluralidade sociocultural dos povos indígenas de Mato Grosso.
A distribuição das \textbf{cargas} \textbf{horárias} do Itinerário \textbf{Formativo}, deve estar articulada com a Formação Geral, de modo a satisfazer as necessidades básicas de aprendizagem dos estudantes.
Ressalta -se que para cada área do conhecimento, são definidas competências específicas, articuladas às respectivas competências atingidas no Ensino Fundamental.
Essas competências específicas de área do Ensino Médio também devem orientar a proposição e o detalhamento das trilhas de aprofundamento a constar no Itinerário \textbf{Formativo}.
O calendário escolar deve ser elaborado de tal maneira que permita que o estudante possa participar das atividades cotidianas, rituais e socioculturais da comunidade : roca, caça, pesca, festas, jogos, reuniões e ritos de passagem.
Em situação específica podem ser usados sábados, feriados e períodos de férias previstos no calendário.
O calendário da escola deve considerar as atividades de acordo com as práticas culturais de cada comunidade, respeitando as especificidades de cada povo, contemplando o que dispõe a lei sobre a quantidade de dias letivos e \textbf{carga} \textbf{horária} estabelecida.
Tais atividades ganham eficácia e serão considerados dias letivos, quando estiverem devidamente registrados no Projeto Político Pedagógico, no Regimento Escolar e no calendário escolar.
A Educação Escolar Indígena deve contribuir para o projeto \textbf{societário} e para o bem viver de cada comunidade indígena, contemplando ações voltadas à manutenção e preservação de seus territórios e dos recursos neles existentes.
O \textbf{currículo} das escolas indígenas, ligado às concepções e práticas que definem o papel sociocultural da escola, diz respeito aos modos de organização dos tempos e espaços da escola, de suas atividades pedagógicas, das relações sociais tecidas no cotidiano escolar, das interações do ambiente educacional com a sociedade, das relações de poder presentes no fazer educativo e nas formas de conceber e construir conhecimentos escolares, constituindo parte importante dos processos sociopolíticos e culturais de construção de identidades.
Nessa perspectiva, os \textbf{currículos} da Educação Básica na Educação Escolar Indígena, como expresso anteriormente, devem \textbf{atender} a perspectiva intercultural e devem ser construídos a partir dos valores e interesses etnopolíticos das comunidades indígenas em relação aos seus projetos de sociedade e de escola.
As conversações e as narrativas estão muito presentes no processo de execução do \textbf{currículo} escolar indígena, assim, compensa refletir sobre, como expressa Ferraço e Magalhães ( 2012 ), potencializar a \textbf{política} das conversações e/ou narrativas, como forma de capacitar os indivíduos e os grupos que se coloquem para produzir e trocar conhecimentos, gerando formas e forças comunitárias, com vista a melhorar os processos de aprendizagem e criação nas coletividades locais.
Nessa mesma perspectiva, Santos ( 2008 ) expressa que as conversações e as narrativas na constituição dos territórios curriculares, para além de sua concepção como grades curriculares, buscam a integração curricular, onde as teorias e os componentes curriculares não são rivais, mas são capazes de integrar os objetivos para a finalidade da aprendizagem.
Com isso, valoriza -se os interesses dos estudantes no que diz respeito à compreensão dos fundamentos pedagógicos da integração curricular.
Para reflexão sobre a integração curricular, trouxemos as contribuições de Beane, a partir de Ferraço e Magalhães ( 2012 ) : > A integração curricular como uma teoria da concepção curricular que está preocupada em aumentar as possibilidades para a integração pessoal, e social através da organização de um \textbf{currículo} em torno de problemas e de questões significativas, identificadas em conjunto por educadores e jovens, independentemente das linhas de demarcação das disciplinas ( BEANE apud FERRAÇO e MAGALHÃES, 2012, p. 224 ).
De modo geral, o \textbf{currículo} integrado deve : estar organizado a partir de questões que tenham significado pessoal e social em situações cotidianas ; valorizar as experiências de aprendizagens que foram significativas ; promover uma formação que priorize valores relativos ao bem comum ; favorecer os conhecimentos relevantes para a sociedade mais ampla e não apenas os de interesse das elites ; e, finalmente, deve estar imbuído de uma concepção de integração, para além de apenas uma \textbf{técnica} alternativa à organização disciplinar.
A integração curricular está relacionada com os problemas reais do cotidiano do estudante e, de acordo com a BNCC, DRC-MT, os objetos do conhecimento, que lhe são pertinentes não devem ser trabalhados de forma casual e pontual, nem como projetos e pesquisas isoladas, visto que devem constituir a base do processo educativo.
Nesse contexto, os problemas que envolvem a comunidade tanto escolar quanto indígena devem ser transformados em temáticas vinculadas ao cotidiano da comunidade escolar que têm despertado na sociedade o anseio por questões de inclusão social, os quais serão transformadas em ação concreta que tornaram o \textbf{currículo} real.
Ao elaborar o Projeto Político Pedagógico, a escola promove uma formação humana capaz de valorar as ações coletivas e o autoconhecimento do estudante, conseguindo dirimir dúvidas e oportunizar o protagonismo indígena e a elaboração do seu projeto de vida, no qual estejam imbuídos o respeito e a tolerância para a construção do saber.
No que \textbf{respalda} os Direitos Humanos : a educação é compreendida como um direito em si mesmo e um meio indispensável para o acesso a outros direitos.
O PPP deve explicitar as ligações entre os diferentes componentes curriculares de forma relacional, assim como, fazer a conexão com situações vivenciadas pelos estudantes, contribuindo para trazer contexto e contemporaneidade à formação das competências e habilidades apresentadas na BNCC e aos objetos de conhecimento, descritos na DRC-MT.
Para o alcance dos propósitos aqui apresentados, as metodologias ativas são uma possibilidade de alcançar as habilidades e competências estabelecidas para esta etapa da Educação Básica.
Principalmente devido ao seu caráter colaborativo, em que as soluções de problemas reais têm potencial pedagógico.
As metodologias ativas fundamentam -se em estratégias de ensino baseadas nas concepções pedagógicas reflexivas e críticas, onde se pode interpretar e intervir sobre a realidade das Aldeias, promover a interação entre as pessoas e valorizar a construção do conhecimento, os saberes e as situações de aprendizagem.
O princípio associado ao uso das metodologias ativas consiste em deslocar o eixo principal da responsabilidade pelo processo de aprendizagem do professor para o estudante.
Nesse sentido, a proposta procura apresentar situações de ensino que despertem o senso crítico do estudante com a realidade, que o faça refletir sobre problemas desafiadores, identifique e organize hipóteses de soluções que mais se adequem e as aplique.
Os princípios que constituem as metodologias ativas são : estudante como centro do ensino e aprendizagem ; Autonomia ; Reflexão ; Problematização da realidade ; Trabalho em equipe ; Inovação e Professor como mediador e/ou facilitador[11 ].
A proposta de trabalho baseada na solução de problemas tem o propósito de tornar o estudante indígena capaz de construir o aprendizado conceitual, procedimental e atitudinal por meio de problemas propostos.
Deste modo, ao expor situações motivadoras, anseia -se minimizar e/ou superar as injustiças sociais e violências de diversas ordens, na perspectiva de compreensão do papel da escola enquanto instituição que edifica valores.
Assim, cabe as Escolas Indígenas trabalhar através de práticas e metodologias inovadoras no sentido de trazer para o \textbf{currículo} temas que envolvem valores culturais, econômicos, éticos e estéticos, valorizando a história do Povo.
Nesse princípio, a Escola Indígena específica e diferenciada, intercultural e bilíngue, \textbf{conquista} a autonomia sociocultural de cada povo ao oportunizar a contextualização na recuperação de sua memória histórica, a reafirmação de sua identidade étnica, o estudo e a valorização da própria língua e da própria ciência – sintetizada em seus etnoconhecimentos, bem como no acesso à informações e aos conhecimentos \textbf{técnicos} e científicos da sociedade majoritária e das demais sociedades, indígenas e não indígenas.
Tendo em vista o conjunto de situações e atividades educativas que os estudantes podem escolher conforme seu interesse, para aprofundar e ampliar aprendizagens em uma ou mais Áreas de Conhecimento e/ou na Formação \textbf{Técnica} e Profissional, as Trilhas de Aprofundamento podem ser estruturadas com foco em uma área do conhecimento, na formação \textbf{técnica} e \textbf{profissional} ou, também, na mobilização de competências e habilidades de diferentes áreas ( integração entre áreas ).
Os quatro eixos estruturantes[12 ] para os \textbf{Itinerários} \textbf{Formativos} devem compor e estar integrados nesta proposta. [ 8 ] : Ensino por etapas se dá quando há necessidade de reunião dos estudantes em um local, tendo em vista que os mesmos vivem dispersos. [ 9 ] : Aprendizagem em construção diz respeito a uma maneira de avaliação que permite \textbf{atender} povos que realizam migração sazonal, ou que vivem em constantes rituais.
Estando matriculados não necessitam apresentar resultados finais no ano letivo civil. [ 10 ] : Mediante consulta livre, prévia e informada conforme Convenção N{\EmojiFont °} 169 da OIT. [ 11 ] : O professor facilitador é uma pessoa que assiste um grupo com objetivos em comum, inclusive ajudando a criá -los, além de alcançá -los.
Ele deve se abster de controlar processos sempre que possível, intervindo apenas para corrigir técnicas ou apresentar questões norteadoras que possam indicar novos caminhos.

% matched lemmas: adolescente, alternância, atender, carga, conquistar, currículo, curso, diretriz, eletivo, entidade, estadual, flexível, formativo, frequentar, horário, itinerário, legislação, oferta, ofertar, orientador, política, profissional, prosseguimento, qualidade, regime, respaldar, societário, técnico
\end{textsample}
